\scanpage{92}
% archive sheet number: 91
Ceci dit, on doit pouvoir \inserted{sans mal} vérifier ceci

\uline{Prop} l'antipodisme $\underline{a}$ est la symétrie \unsure{par r. au|par au} pt $o$ donné par (15)\inserted{si $2$ inv.}, et dans tous les cas il est donné par la formule
\[
\underline{a}(x,y,z)=(\alpha'-x,(1-\alpha')-y,(1-\alpha')-z)\tag*{(17)}
\]
\[
\deleted{\unsure{=(\alpha'-x,-(1+\alpha)-y,\alpha+2-z)}}
\]
Si $2$ est inv., \deleted{$\unsure{a}$}\inserted{$\underline{a}$} a $o$ comme unique pt fixe -- or les $s_i,\deleted{\unsure{s_3}},s'_3,s''_3$ sont $\ne o$ (n'étant pas fixes sous $G$) En car. \deleted{\unsure{zéro}} $2$, $\underline{a}$ est la translation par le vecteur $(\alpha',1,1)$, donc n'a pas de pt fixe (à distance finie).

\uline{NB} En car. $2$, \unsure{\deleted{les solutions}|\uline{les solutions}} de l'équation $\alpha^2+\alpha-1=0$ sur une ext. \inserted{alg. clos. $K$} de $\mathbb{F}_2$ sont les deux éléments de $\deleted{\unsure{\mathbb{F}_4}}\unsure{\mathbb{F}_4}.\ \unsure{\widetilde{K}\setminus K}$, où $\unsure{K|\widetilde{K}}$ est l'unique ext. quadratique de $\mathbb{F}_2$ dans $K$. \unsure{Si on pose}
\[
\mathbb{F}_4\overset{\unsure{\text{déf}|\text{dfn}}}{=}\mathbb{F}_2[T]/(T^2+T+1)=\begin{gathered}\text{corps des racines}\\\text{primitives }3^{\text{èmes}}\text{ de}\\\text{l'unité sur }\unsure{\mathbb{F}_2|\mathbb{F}_p}\end{gathered}
\]
la donnée d'un $\alpha$ sur $S/\mathbb{F}_2$ \unsure{satisfaisant} l'équation $\alpha^2+\alpha-1=0$ équivaut à la donnée d'une $\mathbb{F}_4$-structure sur $S$, i.e. d'un morphisme $S\longrightarrow\operatorname{Spec}\mathbb{F}_4$. Notons que $\alpha^2=\frac1\alpha=\alpha+1$ est l'\uline{autre} él. de $\mathbb{F}_4\setminus\mathbb{F}_2$.
